% Part: incompleteness
% Chapter: incompleteness-provability
% Section: lob-thm

\documentclass[../../../include/open-logic-section]{subfiles}

\begin{document}

\olfileid[id]{inc}{inp}{lob}

\olsection{Teorema L\"ob}

Kalimat G\"odel bagi suatu teori~$\Th{T}$ adalah titik tetap dari
$\lnot \OProv[\Th{T}](y)$, yakni !!a{sentence}~$!G$ sedemikian sehingga
\[
\Th{T} \Proves \lnot \OProv[\Th{T}](\gn{!G}) \liff !G.
\]
Kalimat itu tidak !!{derivable}, karena jika $\Th{T} \Proves !G$, maka
(a)~berdasarkan syarat !!{derivability}~(1),
$\Th{T} \Proves \OProv[\Th{T}](\gn{!G})$, dan (b)~$\Th{T} \Proves !G$
bersama dengan
$\Th{T} \Proves \lnot \OProv[\Th{T}](\gn{!G}) \liff !G$ menghasilkan
$\Th{T} \Proves \lnot \OProv[\Th{T}](\gn{!G})$, sehingga $\Th{T}$ akan
tak konsisten. Sekarang wajar untuk menanyakan status titik tetap dari
$\OProv[\Th{T}](y)$, yakni !!a{sentence}~$!H$ sedemikian sehingga
\[
\Th{T} \Proves \OProv[\Th{T}](\gn{!H}) \liff !H.
\]
Jika kalimat itu !!{derivable}, maka berdasarkan syarat~(1),
$\Th{T} \Proves \OProv[\Th{T}](\gn{!H})$; tetapi kesimpulan yang sama
diperoleh jika kita menerapkan modus ponens pada ekuivalensi di atas. Dengan
demikian, kita tidak memperoleh bahwa $\Th{T}$ tak konsisten, setidaknya
tidak melalui argumen yang sama seperti dalam kasus kalimat G\"odel. Tentu saja,
hal ini tidak menunjukkan bahwa $\Th{T}$ \emph{memang} !!{derive}~$!H$.

Kita dapat memperoleh kemajuan atas pertanyaan ini jika kita membuatnya
sedikit lebih umum. Arah dari kiri ke kanan pada ekuivalensi titik tetap,
$\OProv[\Th{T}](\gn{!H}) \lif !H$, merupakan contoh suatu skema umum yang
disebut \emph{prinsip refleksi}: $\OProv[\Th{T}](\gn{!A}) \lif !A$.
Disebut demikian karena dalam suatu arti prinsip itu menyatakan bahwa
$\Th{T}$ dapat ``merefleksikan'' hal-hal yang teori tersebut dapat
!!{derive}; pada dasarnya, prinsip itu mengatakan, ``Jika $\Th{T}$ dapat
!!{derive}~$!A$, maka~$!A$ benar,'' untuk sebarang~$!A$. Tentu saja, hal ini
hanya benar bagi teori yang sahih, dan ini menunjukkan bahwa pada umumnya
teori tidak akan !!{derive} setiap contohnya. Lalu, contoh-contoh mana yang
suatu teori (yang cukup kuat dan memenuhi syarat
!!{derivability}) dapat !!{derive}? Tentu saja semua contoh dengan $!A$
itu sendiri !!{derivable}. Hanya itulah contoh-contohnya, sebagaimana ditunjukkan oleh hasil
berikut.

\begin{thm}\ollabel{thm:lob}
Misalkan $\Th{T}$ adalah suatu teori yang !!a{axiomatizable} dan memperluas
$\Th{Q}$, dan andaikan $\OProv[\Th{T}](y)$ adalah formula yang memenuhi
syarat P1--P3 dari \olref[2in]{sec}. Jika $\Th{T}$ !!{derive}s
$\OProv[\Th{T}](\gn{!A}) \lif !A$, maka sesungguhnya $\Th{T}$ !!{derive}s
$!A$.
\end{thm}

Dengan kata lain, jika $\Th{T} \Proves/ !A$, maka $\Th{T} \Proves/
\OProv[\Th{T}](\gn{!A}) \lif !A$. Hasil ini dikenal sebagai teorema L\"ob.

\begin{explain}
Heuristik bagi pembuktian teorema L\"ob adalah sebuah pembuktian cerdik
bahwa Sinterklas ada. (Jika Anda tidak menyukai kesimpulan itu, Anda bebas
menggantinya dengan kesimpulan lain apa pun yang Anda inginkan.) Berikut
argumennya:
\begin{enumerate}
\item Misalkan $X$ adalah kalimat, ``Jika $X$ benar, maka Sinterklas ada.''
\item Andaikan $X$ benar.
\item Maka apa yang dikatakannya berlaku; yakni, kita memperoleh: jika $X$
  benar, maka Sinterklas ada.
\item Karena kita mengasumsikan $X$ benar, berdasarkan modus ponens dari
  (2) dan~(3), kita dapat menyimpulkan bahwa Sinterklas ada.
\item Kita telah berhasil menurunkan (4), ``Sinterklas ada,'' dari
  asumsi~(2), ``$X$ benar.'' Berdasarkan pembuktian bersyarat, kita telah
  menunjukkan: ``Jika $X$ benar, maka Sinterklas ada.''
\item Namun, itu tidak lain adalah kalimat~$X$. Jadi, kita telah menunjukkan
  bahwa $X$ benar.
\item Namun kemudian, berdasarkan argumen (2)--(4) di atas, Sinterklas ada.
\end{enumerate}
Formalisasi gagasan ini, dengan mengganti ``benar'' dengan
``!!{derivable},'' dan ``Sinterklas ada'' dengan~$!A$, menghasilkan
pembuktian teorema L\"ob. Siasatnya ialah menerapkan lema titik tetap pada
!!{formula}~$\OProv[\Th{T}](y) \lif !A$. Titik tetapnya bersesuaian dengan
kalimat~$X$ dalam sketsa di atas.
\end{explain}

\begin{proof}[Pembuktian \olref{thm:lob}]
Andaikan $!A$ adalah !!a{sentence} sedemikian sehingga $\Th{T}$ !!{derive}s
$\OProv[\Th{T}](\gn{!A}) \lif !A$. Misalkan $!B(y)$ adalah
!!{formula}~$\OProv[\Th{T}](y) \lif !A$, dan gunakan lema titik tetap untuk
menemukan !!a{sentence}~$!D$ sedemikian sehingga $\Th{T}$ !!{derive}s
$!D \liff !B(\gn{!D})$. Maka setiap pernyataan berikut !!{derivable}
dalam~$\Th{T}$:
\begin{align}
  & !D \liff (\OProv[\Th{T}](\gn{!D}) \lif !A) \ollabel{L-1}\\
  & \qquad \text{$!D$ adalah titik tetap dari~$!B(y)$}\notag \\
  & !D \lif (\OProv[\Th{T}](\gn{!D}) \lif !A) \ollabel{L-2}\\
  & \qquad\text{dari \olref{L-1}}\notag\\
  & \OProv[\Th{T}](\gn{!D \lif (\OProv[\Th{T}](\gn{!D}) \lif !A)}) \ollabel{L-3}\\
  & \qquad \text{dari \olref{L-2} berdasarkan syarat P1}\notag \\
  & \OProv[\Th{T}](\gn{!D}) \lif \OProv[\Th{T}](\gn{\OProv[\Th{T}](\gn{!D}) \lif !A})
  \ollabel{L-4}\\
  &\qquad \text{dari \olref{L-3} dengan menggunakan syarat P2}\notag \\
  & \OProv[\Th{T}](\gn{!D}) \lif (\OProv[\Th{T}](\gn{\OProv[\Th{T}](\gn{!D})}) \lif \OProv[\Th{T}](\gn{!A})) \ollabel{L-5}\\
  &\qquad \text{dari \olref{L-4} dengan menggunakan P2 sekali lagi} \notag\\
& \OProv[\Th{T}](\gn{!D}) \lif \OProv[\Th{T}](\gn{\OProv[\Th{T}](\gn{!D})}) \ollabel{L-6}\\
  & \qquad\text{berdasarkan syarat !!{derivability} P3} \notag\\
  & \OProv[\Th{T}](\gn{!D}) \lif \OProv[\Th{T}](\gn{!A}) \ollabel{L-7} \\
  &\qquad\text{dari \olref{L-5} dan \olref{L-6}}\notag\\
  & \OProv[\Th{T}](\gn{!A}) \lif !A \ollabel{L-8}\\
  &\qquad\text{berdasarkan asumsi teorema} \notag\\
  & \OProv[\Th{T}](\gn{!D}) \lif !A \ollabel{L-9}\\
  &\qquad\text{dari \olref{L-7} dan \olref{L-8}}\notag\\
  & (\OProv[\Th{T}](\gn{!D}) \lif !A) \lif !D \ollabel{L-10}\\
  & \qquad \text{dari \olref{L-1}}\notag \\
  & !D \ollabel{L-11}\\
  & \qquad\text{dari \olref{L-9} dan \olref{L-10}}\notag \\
  & \OProv[\Th{T}](\gn{!D}) \ollabel{L-12}\\
  & \qquad\text{dari \olref{L-11} berdasarkan syarat~P1}\notag \\
  & !A \qquad\qquad\text{dari \olref{L-9} dan \olref{L-12}}\notag
\end{align}
\end{proof}

Dengan teorema L\"ob, tersedia pembuktian singkat bagi teorema
ketaklengkapan kedua (untuk teori dengan predikat !!{derivability} yang
memenuhi syarat P1--P3): jika
$\Th{T} \Proves \OProv[\Th{T}](\gn{\lfalse}) \lif \lfalse$, maka
$\Th{T} \Proves \lfalse$. Jika $\Th{T}$ konsisten,
$\Th{T} \Proves/ \lfalse$. Jadi,
$\Th{T} \Proves/ \OProv[\Th{T}](\gn{\lfalse}) \lif \lfalse$, yakni
$\Th{T} \Proves/ \OCon[\Th{T}]$. Kita juga dapat menerapkannya untuk
menunjukkan bahwa~$!H$, titik tetap dari $\OProv[\Th{T}](x)$,
!!{derivable}. Sebab, karena
\begin{align*}
  \Th{T} & \Proves \OProv[\Th{T}](\gn{!H}) \liff !H\\
  \intertext{maka khususnya}
    \Th{T} & \Proves \OProv[\Th{T}](\gn{!H}) \lif !H
\end{align*}
dan dengan demikian, berdasarkan teorema L\"ob, $\Th{T} \Proves !H$.

% Going in the other direction, for homework I
% may ask you to work through a short proof of L\"ob's theorem, using
% the second incompleteness theorem instead of the fixed-point lemma.

\begin{prob}
Misalkan $\Th{T}$ suatu teori yang diaksiomatisasi secara komputabel, dan
misalkan $\OProv[\Th{T}]$ adalah predikat !!{derivability} bagi $\Th{T}$.
Perhatikan empat pernyataan berikut:
\begin{enumerate}
\item Jika $T \Proves !A$, maka $T \Proves \OProv[\Th{T}](\gn{!A})$.
\item $T \Proves !A \lif \OProv[\Th{T}](\gn{!A})$.
\item Jika $T \Proves \OProv[\Th{T}](\gn{!A})$, maka $T \Proves !A$.
\item $T \Proves \OProv[\Th{T}](\gn{!A}) \lif !A$
\end{enumerate}
Dalam kondisi apa masing-masing pernyataan ini benar?
\end{prob}

\end{document}
